\chapter*{Introduction générale}
\addcontentsline{toc}{chapter}{Introduction générale}
\markboth{Introduction générale}{}

Dans le cadre de la formation d'ingénieur en Informatique et Réseaux à l'École Marocaine
des Sciences de l'Ingénieur (EMSI), les étudiants doivent effectuer un stage de
perfectionnement en entreprise afin de mettre en pratique leurs connaissances théoriques
dans un contexte professionnel réel. C'est dans ce cadre que ce stage de deux mois s'est
déroulé au sein du département Informatique du Mövenpick Hôtel de Tanger.

Le réseau local d'un hôtel comme le Mövenpick voit se connecter, au fil du temps, un
grand nombre d'équipements -- postes fixes, imprimantes, terminaux mobiles, objets
connectés -- sans que le département IT dispose toujours d'un moyen simple d'en garder
un inventaire à jour. C'est dans ce contexte qu'est né le projet \textbf{NetSentry} :
offrir au département IT un outil interne, simple et autonome, permettant de découvrir,
cartographier et surveiller les équipements connectés au réseau local.

Le présent document constitue le \textbf{rapport de stage} du projet NetSentry. Il a
pour objectif de présenter, de manière claire et structurée, le contexte du stage, la
méthodologie de travail suivie, la modélisation et la réalisation technique du système
développé, ainsi que les résultats obtenus.

Ce rapport s'organise en quatre chapitres.

Le premier présente l'entreprise d'accueil, Mövenpick Hôtel de Tanger, son secteur
d'activité, son historique, son organisation interne ainsi que son environnement
(clientèle, fournisseurs, concurrence), puis le département Informatique (IT) au sein
duquel le stage s'est déroulé et la mission confiée durant ces deux mois.

Le deuxième décrit la méthodologie de travail suivie pour la réalisation du projet :
l'analyse du besoin (exigences fonctionnelles et non fonctionnelles), le cahier des
charges qui a cadré le stage, ainsi que les outils et technologies retenus.

Le troisième détaille la modélisation du système NetSentry -- architecture réseau,
diagramme de cas d'utilisation, diagramme de séquence et modèle conceptuel de données --
ainsi que l'organisation du travail retenue pour la réalisation du projet (planning
prévisionnel et suivi des tâches).

Le quatrième présente, à travers des captures d'écran, le fonctionnement réel de
l'outil développé : tableau de bord, détail d'un équipement, alertes, historique des
scans, ainsi que l'historique de versionnage du projet sur GitHub.
